\section{Results: Validation and Comparative Analysis of Causal Discovery Methods}

\subsection{Method Validation on fMRI Network Data}

\subsubsection{GCM Analysis Replicates and Extends Published Network Switching Findings}

To establish the validity of causal discovery approaches for inferring large-scale brain network relationships, we first implemented Granger Causality Mapping (GCM; Roebroeck et al., 2005) and applied it to resting-state fMRI data from 310 subjects in the fBIRN dataset. Following Sridharan et al. (2008), we focused on six ICA-derived brain regions spanning three canonical functional networks: the Central Executive Network (CEN; rPPC, rDLPFC), the Salience Network (SN; rFIC, ACC), and the Default Mode Network (DMN; PCC, VMPFC). 

Our GCM analysis revealed a pattern of directed Granger-causal influences that showed substantial agreement with the key findings of Sridharan et al. (Figure~\ref{fig:gcm_results}A-B). First, we confirmed the central role of the right fronto-insular cortex (rFIC) in network interactions, with rFIC exhibiting significant Granger-causal relationships to regions in both the CEN and DMN. The most prominent finding was robust bidirectional coupling between rFIC and rDLPFC, present in 58\% of subjects—by far the strongest connection in our analysis. This replicates Sridharan et al.'s emphasis on rFIC-DLPFC connectivity while revealing its reciprocal nature in resting-state data. Second, consistent with the proposed role of rFIC in engaging the DMN, we observed directed influences from rFIC to both VMPFC (34\% of subjects) and PCC (30\% of subjects). Third, the anterior cingulate cortex (ACC), identified by Sridharan et al. as a critical partner with rFIC in the salience network, showed directed influences to DMN regions (ACC $\rightarrow$ PCC: 37.4\%, ACC $\rightarrow$ VMPFC: 36.5\%), supporting the salience network's role in regulating default mode activity.

An important extension emerging from our large-sample analysis ($N=310$ vs. $N=18-25$ in Sridharan et al.) was the bidirectional nature of the rFIC-rDLPFC connection. While Sridharan et al. emphasized rFIC $\rightarrow$ DLPFC as a unidirectional control signal during salient events, we found equally strong reciprocal influence (rDLPFC $\rightarrow$ rFIC: 58\%). This bidirectional pattern likely reflects differences between task-evoked transient responses and sustained resting-state coupling, suggesting that while rFIC may initiate network switches during discrete events, salience and executive systems maintain reciprocal communication during ongoing cognition. Collectively, these findings validate our GCM implementation against established results while extending our understanding of network dynamics at rest.

\begin{figure}[htbp]
    \centering
    \includegraphics[width=0.9\textwidth]{figures/gcm_combined_figure.png} % Placeholder
    \caption{\textbf{GCM Analysis Validates Key Network Switching Findings.} \textbf{A.} Network diagram showing Granger-causal relationships between six brain regions from 310 subjects (fBIRN dataset) at 30\% frequency threshold. Node colors indicate functional networks: red = Central Executive Network (CEN: rPPC, rDLPFC), blue = Salience Network (rFIC, ACC), green = Default Mode Network (DMN: PCC, VMPFC). Edge thickness and opacity scale with frequency across subjects. The dominant bidirectional coupling between rFIC and rDLPFC (58\%, thick black edges) replicates the rFIC-DLPFC connectivity emphasized by Sridharan et al. (2008) while revealing its reciprocal nature in resting-state data. Multiple connections between salience and both CEN and DMN regions confirm rFIC's role in mediating network interactions. \textbf{B.} Quantitative ranking of the 15 strongest Granger-causal connections. Bar colors indicate connection types: purple = Salience$\rightarrow$CEN, orange = Salience$\rightarrow$DMN, red = CEN connections, green = DMN connections. The bidirectional rFIC $\leftrightarrow$ rDLPFC relationship dominates (58.06\%), followed by within-DMN connections (PCC $\leftrightarrow$ VMPFC: $\sim$39\%) and salience-to-DMN pathways (ACC $\rightarrow$ PCC/VMPFC: 36-37\%).}
    \label{fig:gcm_results}
\end{figure}

\subsubsection{RASL Reveals Complementary Structural Causal Architecture}

Having validated our approach with GCM, we next applied our Restricted Anytime Structure Learner (RASL) to the same data. RASL employs constraint-based causal structure learning to infer directed acyclic graphs representing structural causal relationships, as opposed to GCM's temporal precedence relationships. This distinction is critical: while GCM asks ``which region's activity temporally precedes another's?'' (Granger causality), RASL asks ``which region structurally causes another in the sense that interventions would propagate?'' (Pearl causality). These represent different but complementary notions of causality.

The RASL analysis identified a markedly different primary causal relationship (Figure~\ref{fig:rasl_results}A-B). The strongest and most consistent structural connection was VMPFC $\rightarrow$ rFIC, present in 89.7\% of the 2,496 causal graph solutions identified across all subjects. This finding was remarkably robust: at a stringent 70\% threshold, VMPFC $\rightarrow$ rFIC was the sole surviving edge, and it remained dominant even at 80\% threshold. Additional strong structural pathways included rDLPFC $\rightarrow$ ACC (67.3\%), PCC $\rightarrow$ rFIC (64.1\%), and VMPFC $\rightarrow$ PCC (62.6\%), collectively suggesting that VMPFC occupies a position of high causal centrality with structural influences propagating to multiple target networks.

\begin{figure}[htbp]
    \centering
    \includegraphics[width=0.9\textwidth]{figures/rasl_combined_figure.png} % Placeholder
    \caption{\textbf{RASL Identifies VMPFC as Primary Structural Driver.} \textbf{A.} Causal structure learning results from RASL applied to the same 310 subjects and brain regions, displayed at 50\% frequency threshold. Network layout and color scheme match Figure 1 for direct comparison. At this moderate threshold, 28 directed edges are present, showing extensive interconnectivity across all three functional networks. The thickest edge (VMPFC $\rightarrow$ rFIC) indicates the dominant structural pathway. \textbf{B.} Top 15 strongest structural causal connections ranked by frequency across 2,496 candidate graph solutions (310 subjects $\times$ $\sim$8 solutions per subject). The VMPFC $\rightarrow$ rFIC connection (89.7\%, green bar) far exceeds all others, emerging as the most consistent structural pathway in the dataset. This finding suggests that the DMN hub (VMPFC) occupies a position upstream of the salience network (rFIC) in the brain's structural causal architecture. Strong connections from VMPFC to multiple target regions (rFIC, PCC, rPPC, rDLPFC) indicate VMPFC's role as a central causal driver across networks.}
    \label{fig:rasl_results}
\end{figure}

\subsubsection{Methodological Complementarity: Temporal Dynamics vs. Structural Constraints}

Direct comparison of GCM and RASL results (Figure~\ref{fig:method_comparison}) revealed both convergence and divergence, illuminating the complementary nature of these approaches. Both methods identified: (1) high interconnectivity between all three functional networks, with directed influences crossing network boundaries; (2) central positioning of the salience network, with connections to both CEN and DMN regions; and (3) active causal roles for DMN regions, challenging purely passive conceptualizations of the default mode network.

The critical difference lay in the identity of the primary network hub and the directionality of the dominant connection. GCM identified \textbf{rFIC $\leftrightarrow$ rDLPFC bidirectional coupling} (58\%) as the strongest temporal relationship, emphasizing reciprocal information flow between salience monitoring and executive control. In contrast, RASL identified \textbf{VMPFC $\rightarrow$ rFIC unidirectional causation} (89.7\%) as the dominant structural pathway, positioning the DMN hub upstream of the salience network in the causal hierarchy.

This divergence is not contradictory but rather reflects the different causal questions each method addresses. The strong VMPFC $\rightarrow$ rFIC structural connection revealed by RASL suggests that the brain's architectural constraints position VMPFC to initiate or enable rFIC activity—that is, VMPFC is structurally upstream in terms of how interventions would propagate through the network. Concurrently, the strong rFIC $\leftrightarrow$ rDLPFC temporal coupling revealed by GCM indicates that once engaged, these regions exhibit tight moment-to-moment reciprocal signaling. Together, these findings suggest a multi-level causal organization: structural positioning (VMPFC enables rFIC) constrains which temporal interaction patterns can emerge (rFIC $\leftrightarrow$ rDLPFC reciprocal dynamics), which in turn support functional outcomes (network switching).

\begin{figure}[htbp]
    \centering
    \includegraphics[width=0.9\textwidth]{figures/comparison_figure.png} % Placeholder
    \caption{\textbf{Complementary Insights from Temporal and Structural Causality.} Side-by-side comparison of network organization revealed by GCM (temporal Granger causality) and RASL (structural causality). Left panel shows GCM network at 50\% threshold, emphasizing the bidirectional rFIC $\leftrightarrow$ rDLPFC temporal coupling (thick curved edges). Right panel shows RASL network at 70\% threshold, where only the unidirectional VMPFC $\rightarrow$ rFIC structural pathway survives this stringent criterion. Bottom panels show comparative bar charts with color-coded connection types. \textit{Similarities:} Both methods identify high cross-network connectivity and central salience network positioning. \textit{Key difference:} GCM emphasizes reciprocal Salience-CEN dynamics (rFIC $\leftrightarrow$ rDLPFC: 58\%), while RASL reveals DMN as structural initiator (VMPFC $\rightarrow$ rFIC: 89.7\%). This complementarity demonstrates that temporal precedence relationships (GCM) and structural causal pathways (RASL) represent distinct organizational principles, both essential for understanding brain network coordination.}
    \label{fig:method_comparison}
\end{figure}

Importantly, comparison across clinical groups revealed preservation of core network architecture. The VMPFC $\rightarrow$ rFIC connection was equally robust in both healthy controls (Group 0: 89.7\%, $N=180$) and schizophrenia patients (Group 1: 89.8\%, $N=130$), with no qualitative topological differences at high thresholds (Figure~\ref{fig:group_comparison}). This invariance suggests that the fundamental structural causal organization of large-scale brain networks remains intact in schizophrenia, with potential clinical differences manifesting in connection strength modulation, temporal dynamics, or task-related engagement rather than reorganization of basic causal architecture.

\begin{figure}[htbp]
    \centering
    \includegraphics[width=0.9\textwidth]{figures/group_comparison_figure.png} % Placeholder
    \caption{\textbf{Clinical Invariance of Structural Network Architecture.} Comparison of RASL-derived causal networks between healthy controls (Group 0, $N=180$, left panel) and schizophrenia patients (Group 1, $N=130$, right panel) at 80\% frequency threshold. Both groups exhibit identical network topology: a single dominant connection from VMPFC to rFIC (Group 0: 89.71\%, Group 1: 89.77\%, difference = 0.06\%). This remarkable similarity across clinical populations indicates that the core structural causal architecture of large-scale brain networks is preserved in schizophrenia. The finding suggests that neuropsychiatric differences may arise from altered network dynamics, connection strength modulation, or task-dependent engagement patterns rather than fundamental reorganization of causal pathways. Node arrangement and colors as in Figures 1-3.}
    \label{fig:group_comparison}
\end{figure}

\subsubsection{Implications}

These results make three important contributions. First, they validate both GCM and RASL as viable approaches for inferring causal relationships from fMRI data, with our GCM implementation successfully replicating established findings and our RASL method revealing stable structural patterns across a large sample. Second, they demonstrate that different causal discovery methods access distinct but complementary aspects of brain network organization—temporal dynamics vs. structural constraints—and that both perspectives are necessary for complete understanding. Third, they reveal VMPFC as a previously underappreciated structural driver whose position in the causal hierarchy may enable salience network engagement, extending frameworks that have emphasized rFIC as the primary control node.

More broadly, these findings highlight the importance of methodological pluralism in neuroscience. Just as anatomical connectivity, functional connectivity, and effective connectivity each reveal different organizational principles, temporal and structural causal analyses provide complementary windows into how the brain coordinates activity across distributed systems. The discovery that RASL identifies different dominant connections than GCM—while both remain consistent with the known functional roles of these regions—suggests that careful attention to causal inference methodology and the specific type of causality being assessed can yield novel insights from existing neural data.

% Supplementary Figures (Optional - include in main text if space permits or move to appendix)

\begin{figure}[htbp]
    \centering
    \includegraphics[width=0.9\textwidth]{figures/supp_gcm_thresholds.png} % Placeholder
    \caption{\textbf{Figure S1: GCM Networks Across Multiple Thresholds.} GCM-derived networks at four frequency thresholds: 0\% (all edges), 30\%, 40\%, and 50\%. As threshold increases, the network progressively simplifies, with the bidirectional rFIC $\leftrightarrow$ rDLPFC connection consistently surviving as the most robust temporal relationship. At 50\% threshold, only two edges remain (both directions of rFIC-rDLPFC coupling), demonstrating the dominance of this salience-executive link in temporal dynamics. Color scheme as in Figure 1.}
    \label{fig:supp_gcm_thresholds}
\end{figure}

\begin{figure}[htbp]
    \centering
    \includegraphics[width=0.9\textwidth]{figures/supp_rasl_thresholds.png} % Placeholder
    \caption{\textbf{Figure S2: RASL Networks Across Multiple Thresholds.} RASL-derived networks at five frequency thresholds: 50\%, 60\%, 70\%, 80\%, and 90\%. Progressive thresholding reveals the hierarchical strength of structural connections. At 70\% and 80\% thresholds, only VMPFC $\rightarrow$ rFIC survives, demonstrating its exceptional robustness as a structural pathway. The sharp transition from 28 edges at 50\% to 1 edge at 70\% indicates a clear separation between the dominant VMPFC $\rightarrow$ rFIC pathway and other connections.}
    \label{fig:supp_rasl_thresholds}
\end{figure}

\begin{figure}[htbp]
    \centering
    \includegraphics[width=0.9\textwidth]{figures/supp_connectivity_matrices.png} % Placeholder
    \caption{\textbf{Figure S3: Complete Connectivity Matrices.} Heatmaps showing complete pairwise connectivity for both GCM (left) and RASL (right) methods. Each cell represents the frequency of a directed edge from source (row) to target (column). Diagonal is masked (no self-connections). Color scale: blue (0\%) to red (100\%). Annotations mark connections exceeding 50\% (GCM) or 70\% (RASL). Matrix view reveals that RASL generally shows higher connectivity levels than GCM (mean: 57\% vs 34\%), but both methods identify similar sets of strong connections despite different primary relationships.}
    \label{fig:supp_matrices}
\end{figure}

\begin{table}[htbp]
\centering
\caption{Comparison of key findings between Sridharan et al. (2008) and current GCM implementation}
\label{tab:agreement}
\begin{tabular}{lcc}
\hline
Finding & Sridharan et al. & Current Study \\
\hline
rFIC central role & \checkmark & \checkmark \\
rFIC $\rightarrow$ DLPFC & \checkmark & \checkmark (58\%) \\
rFIC $\rightarrow$ DMN & \checkmark & \checkmark (30-34\%) \\
ACC $\rightarrow$ DMN & \checkmark & \checkmark (36-37\%) \\
Salience mediates CEN-DMN & \checkmark & \checkmark \\
rFIC bidirectional with DLPFC & Not emphasized & \checkmark (58\%) \\
\hline
\end{tabular}
\end{table}

\begin{table}[htbp]
\centering
\caption{Quantitative comparison of GCM and RASL causal discovery methods}
\label{tab:comparison}
\begin{tabular}{lcccc}
\hline
Metric & GCM & RASL & Difference & p-value \\
\hline
Mean edge frequency & 34.02\% & 57.04\% & +23.02\% & $<0.001$ \\
Maximum frequency & 58.06\% & 89.74\% & +31.68\% & N/A \\
Edges $>50\%$ & 2 & 28 & +26 & N/A \\
Strongest connection & rFIC$\leftrightarrow$rDLPFC & VMPFC$\rightarrow$rFIC & -- & -- \\
Correlation (30 edges) & -- & -- & $\rho=0.43$ & $<0.05$ \\
\hline
\end{tabular}
\end{table}

